\section{Síntesis y ampliación}

El hilo principal de esta sesión puede condensarse en nueve conclusiones consecutivas:

\begin{enumerate}
  \item En sus inicios, el campo de la IA estaba fragmentado por características diseñadas a mano, terminología especializada y distintas topology;
  \item el aprendizaje profundo a gran escala unificó primero el lenguaje de la optimización de extremo a extremo;
  \item el context bottleneck fijo de seq2seq impulsó el alignment diferenciable;
  \item el Transformer elevó la attention a etapa principal de comunicación y la integró con las posiciones, las conexiones residuales, la normalización, la MLP, las heads y las masks en un package completo;
  \item Q/K/V pueden entenderse como la necesidad de direccionamiento, la etiqueta que puede emparejarse y el contenido que se transmite;
  \item nanoGPT muestra el ciclo completo de token stream, shifted targets, causal mask, cross-entropy y autoregressive generation;
  \item GPT, BERT y T5 se distinguen principalmente por la connectivity, la state source y el objective;
  \item la reutilización entre modalidades depende de la tokenización, del type/position embedding y del bias propio de cada dominio, y no de simplemente convertir todos los problemas en texto;
  \item la expresividad, la facilidad de optimización, la eficiencia en GPU y la escala producen en conjunto el in-context behavior, lo que permite que el prompt reconfigure el modelo en tiempo de ejecución.
\end{enumerate}

\begin{warningbox}{Límites finales de la evidencia}
La clase demostró cómo implementar un Transformer mínimo y presentó ejemplos de primera mano en varios dominios; no demostró que la attention sea la única arquitectura óptima, ni que el scaling pueda continuar indefinidamente, ni que el in-context learning equivalga al descenso de gradiente, ni que los modelos de lenguaje sean ejecutores confiables de programas arbitrarios. Un aprendizaje de calidad debe marcar por separado el «mecanismo ejecutable», el «fenómeno empírico», el «juicio del ponente» y la «conjetura a futuro».
\end{warningbox}

\subsection{Lecturas adicionales}

Se recomienda leer siguiendo una cadena de preguntas, en lugar de acumular títulos de artículos:

\begin{enumerate}
  \item Bengio et al., \href{https://www.jmlr.org/papers/v3/bengio03a.html}{\emph{A Neural Probabilistic Language Model}}: modelo de lenguaje neuronal de ventana fija y learned embeddings;
  \item Sutskever et al., \href{https://proceedings.neurips.cc/paper/2014/hash/a14ac55a4f27472c5d894ec1c3c743d2-Abstract.html}{\emph{Sequence to Sequence Learning with Neural Networks}}: encoder--decoder de longitud variable;
  \item Bahdanau et al., \href{https://arxiv.org/abs/1409.0473}{\emph{Neural Machine Translation by Jointly Learning to Align and Translate}}: soft alignment y context dinámico;
  \item Vaswani et al., \href{https://arxiv.org/abs/1706.03762}{\emph{Attention Is All You Need}}: el package completo del Transformer;
  \item Karpathy, \href{https://github.com/karpathy/nanoGPT}{\emph{nanoGPT}}: conectar datos, modelo, entrenamiento y generación con un código mínimo;
  \item Dosovitskiy et al., \href{https://arxiv.org/abs/2010.11929}{\emph{An Image is Worth 16x16 Words}}, y Chen et al., \href{https://arxiv.org/abs/2106.01345}{\emph{Decision Transformer}}: comparar cómo distintos dominios definen el token;
  \item Brown et al., \href{https://arxiv.org/abs/2005.14165}{\emph{Language Models are Few-Shot Learners}} y Wei et al., \href{https://arxiv.org/abs/2201.11903}{\emph{Chain-of-Thought Prompting}}: distinguir entre fenómeno de comportamiento, prompt protocol y explicación del mecanismo.
\end{enumerate}

El orden de práctica más eficaz es: primero, verificar línea por línea las tensor shapes y la causal mask sobre Tiny Shakespeare; después, cambiar a un subword tokenizer y a un context más largo; luego, implementar por separado el unmasked encoder y la cross-attention; por último, elegir un dominio no textual y escribir explícitamente la element definition, el position/type encoding, la connectivity y el objective. Solo así el lema «Transformer everywhere» se convierte en un método de diseño verificable.

\end{document}
